\begin{abstract}
% Real-world images often suffer from unknown and intertwined degradations such as noise, blur, and environmental effects. Restoring such images requires not only perceiving the underlying degradations but also handling them adaptively, posing a severe challenge for all-in-one image restoration models. To address this issue, we present Restoramind, an intelligent agent trained end-to-end with reinforcement learning that formulates image restoration as a sequential decision-making process. Specifically, Restoramind first analyzes the input image to identify its degradation patterns, then sequentially selects and applies the most appropriate restoration tools. At each step, the agent re-evaluates intermediate results, reasons about remaining degradations, and refines its restoration plan until the output becomes visually clean. Extensive experiments demonstrate that Restoramind effectively recognizes complex and unknown degradations and progressively restores high-quality images, achieving superior generalization and visual fidelity on Benchmark.
Recent years have witnessed remarkable progress in restoring images with specific degradation types. However, real-world image restoration remains challenging, as natural images often suffer from complex and unknown degradations that pose significant challenges to a single unified model. Inspired by the reasoning capabilities of large multimodal models, recent methods have leveraged them to build intelligent restoration systems, typically following two stages: (1) reasoning about degradation types and generating an initial tool execution sequence, and (2) applying post-search strategies to refine the sequence through trial and error. The iterative nature of this process incurs considerable computational overhead, which hinders its applicability in the wild. In this paper, we present RestoraMind, an LMM-powered agent trained via reinforcement learning to jointly reason about degradations and adaptively utilize restoration tools. Furthermore, we incorporate image statistics as informative priors to guide the agent in selecting the most suitable tool among multiple candidates designed for the same degradation type. Extensive experiments demonstrate that RestoraMind formulates image restoration as a sequential decision-making process, reducing tool invocations to near-linear scale while substantially outperforming previous state-of-the-art methods.
\end{abstract}